\newpage
\section{Introduction}

\subsection{Motivations}

Occupational choice is among the most consequential decisions an individual makes, and it is typically made with limited direct knowledge of the work itself. A prospective entrant selects a field on the basis of job titles, salary figures, and descriptions authored by parties with an interest in recruiting them. The texture of the work --- its daily pressures, its trade-offs, the tasks that occupy most of its hours --- generally becomes apparent only after entry, when the cost of reversing the decision is high.

Research on career decision-making identifies the mechanism behind this difficulty. In the taxonomy of Gati, Krausz and Osipow\cite{r1}, whose structure was recently confirmed across thirteen countries\cite{r2}, insufficient information about occupations constitutes a principal category of decision difficulty, distinct from indecisiveness or inadequate self-knowledge. The obstacle is informational rather than dispositional: individuals are not primarily unable to decide, but unable to decide on an adequate basis.

Conditions in Vietnam illustrate the problem in a particularly compressed form. Entry to higher education is mediated by national examination results and published admission benchmark scores, so the selection of a field of study is effectively completed at seventeen or eighteen, under time pressure, and is constrained by the applicant's score relative to institutional thresholds rather than by any assessment of occupational fit\cite{r3}. The decision is frequently not the applicant's own. A 2026 survey of approximately 9,200 students across ten Vietnamese universities reported that more than one fifth had selected their field on parental direction, that nearly one fifth had entered without a defined occupational goal, and that around a third were subsequently uncertain whether they would choose the same field again\cite{r4}. Labour-market data indicate the downstream cost. Analysing successive waves of the national Labour Force Survey, Tran et al.\cite{r5} found the proportion of graduates in fully matched employment declining over time, with vertical and combined mismatch carrying measurable wage penalties; related analyses associate the same pattern with underutilised human capital at national scale. Employer-reported skills shortages, particularly in science and technology fields, compound the effect.

The prevailing response, in Vietnam and elsewhere, is psychometric. Interest inventories derived from Holland's typology, personality instruments in the Myers-Briggs tradition, and their commercial derivatives dominate career guidance provision. These instruments are informative about preferences but structurally unable to address the difficulty Gati's taxonomy identifies: they elicit what an individual believes about themselves rather than supplying information about occupations, and they cannot establish whether that individual would perform or persist in a given role. Their measurement properties are also contested. Early retest studies found that a substantial share of respondents were assigned a different Myers-Briggs type after only a few weeks\cite{r6}, whereas the publisher reports test--retest coefficients of 0.81 to 0.86 on the underlying continuous scales\cite{r7}; the instability concerns the dichotomised type labels on which guidance is usually based, for respondents whose scores lie near a boundary. The more fundamental limitation is not reliability but what is measured: a stable preference remains a preference, not information about the work.

An alternative is well established in organisational psychology but rarely applied to career exploration. The realistic job preview literature demonstrates that balanced previews, including a role's unattractive features, improve person--job fit and reduce subsequent turnover, with perceived honesty identified as the operative mechanism\cite{r8,r9}. Such previews are, in effect, information about the work delivered before commitment. Their delivery, however, has historically depended on either direct workplace exposure or content authored case by case, neither of which scales to the breadth of occupations an undecided student might wish to consider. Recent advances in generative language models alter this constraint, and that change in feasibility motivates the present project.

\subsection{Objectives}

The overall objective is to enable a person facing an occupational decision to experience the work of an occupation, receive an assessment of how they performed it, and use that assessment to decide which occupations to explore further --- before committing to one. This objective is realised by JobQuest, a deployable platform specified in Chapter 4. The present phase delivers a working system covering the information technology field; the subsequent capstone phase extends its occupational coverage and operational maturity.

The overall objective is pursued through four specific objectives. Each is stated with its scope, a measurable target, and the time by which the target is to be met; the present phase ends with 2026, and the capstone phase follows in 2027.

\begin{enumerate}[leftmargin=1.5em]
  \item To deliver occupational tasks that remain faithful to the occupations they represent, by constraining AI generation with expert-authored material, so that the experience a user receives conveys accurate information about the work. Scope: the information technology roles of the present phase, with three pilot roles for expert evaluation. Target: at least 95\% of the facts a seed requires are preserved in generated scenarios; experts rate generated scenarios with an item-level content validity index of at least .78 and a scale-level average of at least .90; and seeded generation is judged acceptable significantly more often than unseeded generation. Time: fact preservation by the end of the present phase; expert ratings by the end of the capstone phase.
  \item To assess users' responses, including open-ended ones, through automated judgment whose agreement with domain experts is established, so that the result, feedback, and guidance a user receives can be relied upon for formative use. Scope: the free-text response formats in the three pilot roles. Target: quadratic weighted kappa of at least 0.70 between automated and expert scores, a standardised mean score difference of at most 0.15, and agreement no more than 0.10 below that between experts. Time: sensitivity to each rubric criterion demonstrated on the synthetic benchmark by the end of the present phase; agreement with experts by the end of the capstone phase.
  \item To convert demonstrated competence into guidance on which occupations a user is prepared to explore next, so that exploration builds on what the user has shown they can do. Scope: levels within an occupation and adjacency among the roles in the catalogue. Target: a level unlocks at, and only at, its specified competence credit, and the adjacent occupations offered to a set of test profiles match those derived from the occupational taxonomy. Time: the end of the present phase.
  \item To make the platform freely and sustainably available, with occupational content governed by domain experts, so that access to occupational experience does not depend on a user's ability to pay. Scope: all content and all users. Target: no level or occupation is placed behind payment; the mean model cost of a completed session remains within the configured ceiling; and the evidence required to pass the first stage of the roadmap (Section 4.1.5) is in place. Time: the end of the present phase, with later stages gated as the roadmap sets out.
\end{enumerate}

Each specific objective secures one part of the overall objective. The first makes it possible to experience the work of an occupation as it is; the second supplies the assessment of how the user performed it; the third turns that assessment into a decision about which occupations to explore further; and the fourth ensures that this is available before commitment and without regard to means. The first two specific objectives correspond to the research problems stated in Section 4.2.2; the third and fourth are realised through the system's design and its business model (Section 4.1.5).

\subsection{Scope}

The system delivered in this phase is complete in function and bounded in coverage to the information technology field and its specialisations. The coverage boundary is methodological: constructing a reference set of gold-standard scenarios and expert judgments requires domain expertise that the project holds in this field, and without such a set the first two specific objectives cannot be verified. The design is domain-agnostic, so extending coverage to further fields requires expert authoring rather than redesign.

Within this coverage the work comprises the authoring of occupational scenario material by domain experts and their review of the scenarios generated from it; the presentation and publication of that material by content administrators; AI generation and adaptation of task scenarios from that material; automated assessment of structured and open-ended responses, returning a result, feedback, and guidance; a competence-based progression model linking levels and adjacent occupations; the presentation and engagement design described in Section 4.2.1; a supporting orientation conversation for users uncertain where to begin; administration of the occupation catalogue and of the AI models; and an offline evaluation of generation quality and assessment accuracy against the expert reference set.

\subsection{Significance of the Project}

\subsubsection{Significance from the Practical Perspectives}

The practical contribution is a freely accessible means of acquiring occupational information experientially before commitment, for students selecting a field of study, graduates entering the labour market, workers considering transition, and the institutions that guide them. Its value rests on three established findings: that balanced, honest previews improve fit and reduce later regret\cite{r8,r9}; that the obstacle they address, insufficient occupational information, is a principal source of career decision difficulty\cite{r1,r2}; and that performing authentic tasks is the strongest source of self-efficacy, which social cognitive career theory identifies as a determinant of choice and persistence\cite{r10,r11}. Because the platform is designed to be offered without charge and operated in partnership with public education authorities (Section 4.1.5), its reach does not depend on users' ability to pay. This matters particularly in Vietnam, where the decision falls at an age and within a family setting that leave most applicants little opportunity to observe the work they are selecting, and where the existing guidance market consists almost entirely of self-report instruments.

\subsubsection{Significance from the Scientific Perspectives}

The scientific contribution lies in artificial intelligence applied to education, in two problems that current research identifies as open. The first is grounded generation of educational content: language models produce fluent instructional material readily, but fidelity to a specified domain requires explicit constraint, and the project contributes an expert-in-the-loop generation method evaluated against a domain reference set. The second is validated automated assessment: language models are increasingly used in educational technology to judge open-ended learner responses, yet reported agreement is systematically optimistic, and chance-corrected, difficulty-stratified validation is rarely reported. The project contributes such a validation in the setting of behavioural, or stealth, assessment, where the evidence of competence is the learner's performance of a task rather than an answer to a test item.

The application setting is itself a contribution to educational technology. Experiential career education has not previously combined these techniques, and the systems that approach it are research prototypes evaluated through small qualitative or within-subject studies. An evaluation reporting chance-corrected agreement, difficulty stratification, and test--retest consistency offers a more reproducible basis for the claims such systems make.

\subsection{Report Structure}

Chapter 2 establishes the background: the evidence that experiential exposure conveys occupational information more effectively than description, the basis for measuring competence from behaviour, the conditions under which game-based delivery supports rather than displaces learning, relevant career-development theory, the generative and assessment technologies involved, and the data resources and benchmarks available for grounding and evaluation. Chapter 3 reviews related work, examining research prototypes and deployed systems, deriving criteria for comparison, and identifying the gap the project addresses. Chapters 4 to 8 follow the prescribed structure for the specialized project, covering the proposed system, its analysis and design, implementation and testing, evaluation, and conclusions.
